\cotefile{15}
\scanpage{120}\mspage{78}
and since $\Phi_k\ne 0$ in $Q'$, $f_nQ'=Q$, then $\frac{\Phi_{\unsure{0|O}}}{\Phi_k},\ldots,\frac{\Phi_T}{\Phi_k}$, and
with them $F_1$ and $F_2$ are regular in $Q'$; from the relations
$F_2(Q')=F_2(Q)\ne 0$ follows the regularity of $R$ in $Q'$.

Let us now show that, for large $n$, the mapping $f_n$
becomes one-to-one. The equation $\varphi^{\unsure{nn_e|\eta\eta_e|m_e}}=0$ will give a
cross-\unsure{'}section of the variety $V^{(n)}$ by some hyperplane, not
containing $V^{(n)}$. Evidently, this cross-section is the image
of $s-1$-dimensional subvariety $\mathscr E_{\varphi}$ of the variety $V$
(the equation $\mathscr E_{\varphi}$ in $U^{(1)}$ - the set $\varphi^{(i)}=0$).

According to the induction hypothesis, there exists
$m(\mathscr E_{\varphi})$ such that the elements of the group \handwritten{\unsure{$\checkmark$|of}}
$H^{\unsure{0|o}}(\mathscr E_{\varphi},O_{\mathscr E_{\varphi}}[mE']_{\mathscr E_{\varphi}})$ will give a biregular imbedding of
the variety $\mathscr E_{\varphi}$ in a projective space. From the exact
sequence (19) we obtain that this imbedding of the variety
$\mathscr E_{\varphi}$ is the restriction of the mapping $f_m\colon V\to V^{(m)}$ on
the subvariety $\mathscr E_{\varphi}$. We obtain, in particular that the
image $\mathscr E_{\varphi}$ by the mapping $f_m$ possesses a dimension $s-1$.
Hence it follows that the cross-section of the variety $V^{(m)}$
by a hyperplane, not containing $V^{(m)}$, possesses a dimension
$s-1$ and, consequently,
\[
\handwritten{\dim V^{(m)}}=s.
\]
% archive sheet number: 119
